\subsection{模的积}

设 \((\mathbf{E}_i)_{i\in I}\) 是同一环 A 上的一族模。可以立即验证， \(\mathbf{E}_i\) 上的模结构（I，§4，第 8 号）的乘积是乘积集合 \(\mathbf{E} = \prod_{i\in I}\mathbf{E}_i\) 上的一个 A-模结构。带此结构的集合 E 称为诸模 \(\mathbf{E}_i\) 的积模；若 \(x = (x_i)\) , \(y = (y_i)\) 是 E 的两个元素，则

\[
\left\{ \begin{array}{l l} x + y = (x _ {i} + y _ {i}) \\ \lambda . x = (\lambda x _ {i}) & \text { for   all } \lambda \in \mathbf {A}. \end{array} \right.\tag{17}
\]

公式 (17) 表明投影 \(pr_{i}:E\to E_{i}\) 是线性映射；这些映射显然是满射。

回顾，若指标集 I 为空，则乘积集合 \(\prod_{i\in I}E_{i}\) 仅由一个元素组成；此集合上的积模结构即是使这个唯一元素为 0 的那个结构。

命题 4. 设 \(\mathbf{E} = \prod_{i\in I}\mathbf{E}_i\) 是 A-模族 \((\mathbf{E}_i)_{i\in I}\) 的积。对每个 A-模 \(\mathbf{F}\) 与每族线性映射 \(f_i:\mathbf{F}\to \mathbf{E}_i\) ，存在且仅存在一个映射 \(f\) （ \(\mathbf{F}\) 到 \(\mathbf{E}\) 中的映射），使得 \(\mathbf{pr}_i\circ f = f_i\) 对所有 \(i\in I\) 成立，且这个映射是线性的。

这直接由定义得出。

模的积是“结合的”：若 \((J_{\lambda})_{\lambda \in L}\) 是 I 的一个划分，则典范映射

\[
\prod_ {i \in I} E _ {i} \rightarrow \prod_ {\lambda \in L} \left(\prod_ {i \in J _ {\lambda}} E _ {i}\right)
\]

是一个同构。

命题 5. (i) 设 \((\mathbf{E}_t)_{i \in I}\) , \((\mathbf{F}_t)_{i \in I}\) 是具有同一指标集 I 的两个 A-模族；对每族线性映射 \(f_i: \mathbf{E}_t \to \mathbf{F}_t\) ( \(i \in I\) )，映射 \(f: (x_i) \to (f_i(x_i))\) （ \(\prod_t \mathbf{E}_t\) 到 \(\prod_t \mathbf{F}_t\) 中的映射，有时记作 \(\prod_t f_i\) ）是线性的。

\begin{enumerate}
\def\labelenumi{(\roman{enumi})}
\setcounter{enumi}{1}
\tightlist
\item
  设 \((\mathbf{G}_{\iota})_{\iota \in \mathbf{I}}\) 是以 I 为指标集的第三族 A-模，并且对所有 \(t \in I, \text{ 设 } g_t: F_t \to G_t \text{ 是一个线性映射；设 } g = \prod_t g_t. \text{ 对每个序列 } E_t \xrightarrow{f_t} F_t \xrightarrow{g_t} G_t \text{ 为正合，当且仅当序列}\)
\end{enumerate}

\[
\prod_ {i} \mathrm{E} _ {i} \xrightarrow {f} \prod_ {i} \mathrm{F} _ {i} \xrightarrow {g} \prod_ {i} \mathrm{G} _ {i}
\]

是正合的。

断言 (i) 直接由定义得出。另一方面，说 \(y = (y_{i})\) 属于 \(\operatorname{Ker}(g)\) ，就是说 \(g_{i}(y_{i}) = 0\) 对所有 \(i \in I\) 成立，从而就是说 \(y_{i} \in \operatorname{Ker}(g_{i})\) 对所有 \(i \in I\) 成立；类似地，说 \(y\) 属于 \(\operatorname{Im}(f)\) ，就是说存在 \(x = (x_{i}) \in \prod_{i} E_{i}\) 使得 \(y = f(x)\) ，这等价于说 \(y_{i} = f_{i}(x_{i})\) 对所有 \(i \in I\) 成立，或者也说 \(y_{i} \in \operatorname{Im}(f_{i})\) 对所有 \(i \in I\) 成立；由此得 (ii)。

推论。在命题 5 (i) 的条件下，

\[
\operatorname{Ker} (f) = \prod_ {i \in I} \operatorname{Ker} (f _ {i}), \quad \operatorname{Im} (f) = \prod_ {i \in I} \operatorname{Im} (f _ {i})\tag{18}
\]

并且存在典范同构

\[
\operatorname{Coim} (f) \cong \prod_ {i \in I} \operatorname{Coim} (f _ {i}), \quad \operatorname{Coker} (f) = \prod_ {i \in I} \operatorname{Coker} (f _ {i})
\]

它们分别是这样得到的：把元素 \(x = (x_{i})\) （属于 \(\prod_{i} E_{i}\) ，模 \(\operatorname{Ker}(f)\) ）的类（相应地，元素 \(y = (y_{i})\) （属于 \(\prod_{i} F_{i}\) ，模 \(\operatorname{Im}(f)\) ）的类）对应到诸 \(x_{i}\) 模 \(\operatorname{Ker}(f_{i})\) 的类组成的族（相应地，诸 \(y_{i}\) 模 \(\operatorname{Im}(f)\) 的类组成的族）。

特别地，要使 \(f\) 是单射（相应地，满射、双射、零映射），必须且只需对所有 \(\iota \in I\) ， \(f_{\iota}\) 都是单射（相应地，满射、双射、零映射）。

若对所有 \(\iota \in I\) ，我们考虑一个子模 \(F_{\iota}\) （ \(E_{\iota}\) 的子模），则模 \(\prod_{\iota \in I} F_{\iota}\) 是 \(\prod_{\iota \in I} E_{\iota}\) 的一个子模，并且根据命题 5 的推论，存在典范同构

\[
\prod_ {i \in I} \left(E _ {i} / F _ {i}\right)\cong \left(\prod_ {i \in I} E _ {i}\right) / \left(\prod_ {i \in I} F _ {i}\right).\tag{19}
\]

模的积的一个重要例子是所有因子模都等于同一个模 F 的情形；它们的积 \(F^{I}\) 此时正是 I 到 F 的映射的集合。对角映射 \(F \rightarrow F^{I}\) 把 \(x \in F\) 映到 I 上等于 x 的常值函数，它是线性的。若 \((\mathbf{E}_{t})_{t \in I}\) 是一族 A-模，且对所有 \(\iota \in I\) ， \(f_{\iota}: F \to E_{\iota}\) 是线性映射，则线性映射 \(x \mapsto (f_{\iota}(x))\) （ \(F\) 到 \(\prod_{\iota \in I} E_{\iota}\) 中的映射）是映射 \((x_{\iota}) \mapsto (f_{\iota}(x_{\iota}))\) （ \(F^{\mathrm{I}}\) 到 \(\prod_{\iota \in I} E_{\iota}\) ）与对角映射 \(F \to F^{\mathrm{I}}\) 的复合。

